% !PRINTED none
\weberpage{none}

\begin{center}
\gesperrt{Zehnter Abschnitt.}

\bigskip

\gesperrt{Allgemeine Theorie der metacyklischen Gleichungen.}
\end{center}

\bigskip
\begin{center}
\rule{4cm}{0.4pt}
\end{center}
\bigskip

\begin{center}
\S.\,73.

\medskip

\gesperrt{Die Resolventen der Compositionsreihe.}
\end{center}

\bigskip

Aus den S\"atzen der allgemeinen Gruppentheorie, die im ersten Buche dieses Bandes
behandelt sind, ergeben sich wichtige algebraische Folgerungen, wenn man sie auf die
\gesperrt{Galois}'sche Gruppe einer Gleichung anwendet.

Es wird jetzt ein beliebiger K\"orper $\Omega$ als Rationalit\"atsbereich angenommen,
$f(x) = 0$ sei eine Gleichung in diesem K\"orper ohne mehrfache Wurzeln und $P$ ihre
\gesperrt{Galois}'sche Gruppe. Diese Gruppe $P$ habe einen Normaltheiler $Q$. Wir bezeichnen mit
$n$ den Grad von $P$ und mit $j$ den Index des Theilers $Q$, so dass $n : j$ der Grad
von $Q$ ist.

Wir haben im ersten Bande (\S.\,156) gesehen, dass die Gruppe von $f(x)$ durch
Adjunction einer Wurzel einer irreduciblen Normalgleichung $j^{\mathrm{ten}}$ Grades
auf $Q$ reducirt wird und diese H\"ulfsgleichung $j^{\mathrm{ten}}$ Grades haben wir
als \gesperrt{Partialresolvente} bezeichnet.

Die \gesperrt{Galois}'sche Gruppe dieser Partialresolvente haben wir so erhalten: Bedeutet
$\psi$ eine zu der Gruppe $Q$ geh\"orige Function der Wurzeln von $f(x)$ und ist $P$
in die Nebengruppen
\[
P = Q + Qa + Qb + Qc + \cdots
\]
zerlegt, wo also $a, b, c, \ldots$ gewisse Permutationen aus $P$ sind, so geht $\psi$
durch die ganze Nebengruppe $Qa$ in ein und dieselbe Function $\psi_a$ \"uber, und die
\gesperrt{Galois}'sche Gruppe der Resolvente besteht aus den Substitutionen
\[
(\psi, \psi), (\psi, \psi_a), (\psi, \psi_b), (\psi, \psi_c), \ldots
\]
